\section{The aiprov Method}\label{sec:method}

\textbf{Model.} aiprov extends PROV-O with three agent classes
(\texttt{HumanAgent}, \texttt{AIAgent}, \texttt{ToolAgent}), four
activity classes (\texttt{AuthoringPass}, \texttt{AuditPass},
\texttt{Build}, \texttt{Repair}), and entity classes for artefacts,
claims, sources, transcripts, and prompts, where prompts are treated
as source code and bound to activities as plans
(Figure~\ref{fig:vocab}). Humans carry ORCID and affiliation; AI
agents carry model, version, provider, endpoint, context window, and
knowledge cutoff; tool agents cover deterministic actors such as CI
jobs. Treating prompts as plans deserves a word: a prompt is to an AI
activity what a script is to a batch job, the executable
specification of intent, so it is hashed and bound with
\texttt{prov:hadPlan} rather than paraphrased into a comment. An
auditor who wants to know \emph{why} an artefact came out as it did
reads the plan and the transcript; neither is recoverable after the
fact if not captured when the activity runs.

\textbf{Telemetry.} Activities carry session and request identifiers,
token counts by class (input, output, cache, reasoning), sampling
parameters, cost, energy, and the tools invoked. One rule governs all
of it: omit, don't estimate. Absent telemetry is recorded as absent;
invented precision is worse than a visible gap. Generated artefacts
receive sha256 content hashes at logging time, and prompt files
receive prompt hashes.

\textbf{Invariant I: no parentless claim.} Every \texttt{aiprov:Claim}
needs a generating activity, an attributed agent, and a verification
state. The conformance validator treats violations as hard failures
and is suitable as a pre-commit hook or CI gate.

\textbf{Invariant II: the verification ladder.} Rungs answer who
checked what (Figure~\ref{fig:ladder}, Table~\ref{tab:ladder}).
Promotions must strictly climb and are themselves logged as audit
activities, so every rung a claim or source holds has a recorded
grantor. Disagreement is first-class: an agent may \emph{refuse} a
promotion, which changes no state but is logged as an audit activity
carrying the refused rung, so ``checked and not convinced'' is
distinguishable from ``never checked''. Rungs 5--6 are human-only; an AI-granted promotion to them is
a validator hard failure. Two rungs deserve comment. Vendoring
(rung~4) is not a verification statement: an unread vendored copy
proves nothing beyond \texttt{reference-resolved}. Its role is
instrumental, the access gate that makes human verification feasible
and durable. Operationally, the agent must hand the human the
evidence: \texttt{promote} refuses the human-only rungs unless the
source is vendored (path and hash recorded) or carries a clear access
link, and prints the review material with every request; a human
confirmation against link-only material is flagged, since the audit
target remains mutable. At the top, \texttt{human-read} subsumes
\texttt{human-confirmed}: it asserts full-text familiarity \emph{and}
confirmation, because a name alone would mislead, as reading does not
entail checking. Both refinements were operator-contested
mid-experiment (Section~\ref{sec:meta}). The ladder itself has
provenance: it canonicalizes labels field-tested in
\emph{Obscurity-Is-Dead}~\cite{obscurity-is-dead} and hardened in
F(AI)\textsuperscript{2}R~\cite{fair2r-repo}, and legacy names remain
machine-readable aliases, so that ancestor graphs can consolidate
onto this ladder without rewriting history; the aliases are in place,
though the consolidation has not yet been exercised against the
ancestor repositories' graphs.

\textbf{Commit binding.} An activity can only bind a commit that
already exists, and a commit cannot contain the graph entry describing
itself. The working discipline that follows: commit the artefacts, log
the activity with that commit's hash, then commit the graph update
separately. The graph never claims a hash it could not have known; the
cost is one extra commit per logged round. What the two commits buy
is tamper evidence in both directions. The graph entry names a commit
whose tree contains the artefacts it describes, content-hashed; the
graph itself lives in later commits of the same history. Backdating a
record would require rewriting published git history, and quietly
swapping an artefact would break its recorded hash: either
falsification is possible, but neither is silent, and that asymmetry,
honest work is cheap and dishonest work is loud, is the design goal
throughout.

\textbf{The validator, exhaustively.} Conformance is seven concrete
checks, each a SPARQL query over the graph, and naming them makes the
invariants operational rather than aspirational. Three enforce
Invariant~I: no claim without a generating activity, no claim without
an attributed agent, no activity without an associated agent. Two
enforce Invariant~II: no human-only rung granted by an AI agent, and
no human-only rung without a recorded promotion activity naming a
human grantor, which closes the loophole of a rung edited directly
into the graph without an audit trail. Two enforce the access gate:
no source above \texttt{needs-research} without a resolvable DOI, URL,
or vendored copy, and a warning, rather than a failure, for
human-verified sources whose audit target is link-only and therefore
mutable. Six checks are hard failures with a nonzero exit code;
that exit code is the entire integration contract, which is why the
same command serves unchanged as a pre-commit hook, a CI gate, and an
auditor's first step.

\textbf{Auditability operations.} \texttt{report} aggregates tokens,
cost, and rung distributions; \texttt{extract} computes the backward
closure of everything that contributed to chosen entities;
\texttt{dashboard} renders the graph as a self-contained offline
ledger. Sources enter through open registries (Crossref, OpenAlex,
arXiv, DataCite): a search command sweeps them, and registration with
DOI verification promotes a source to \texttt{reference-resolved} only
if the registry actually resolves it.

\begin{figure}[t]
\centering
\resizebox{\columnwidth}{!}{%
\begin{tikzpicture}[
  font=\scriptsize, node distance=2.5mm,
  ecls/.style={draw=orange!60!black, fill=yellow!22, rounded corners=1pt,
               inner sep=2.5pt, font=\tiny\ttfamily, align=center},
  acls/.style={draw=blue!50!black, fill=blue!10, rounded corners=1pt,
               inner sep=2.5pt, font=\tiny\ttfamily, align=center},
  gcls/.style={draw=orange!75!black, fill=orange!25, rounded corners=1pt,
               inner sep=2.5pt, font=\tiny\ttfamily, align=center},
  pkg/.style={draw=black!45, rounded corners=2pt, inner sep=3.5pt},
  pkl/.style={font=\tiny\itshape, text=black!60},
  e/.style={-{Stealth[length=1.7mm]}},
  lbl/.style={font=\tiny\itshape, fill=white, inner sep=1pt}]
% Entities (top)
\node[ecls] (claim) at (0,0)
  {Claim\\[-2pt]\tiny verificationState};
\node[ecls, right=of claim] (artefact)
  {Artefact\\[-2pt]\tiny sha256, commit};
\node[ecls, right=of artefact] (source)
  {Source\\[-2pt]\tiny doi, rung, self};
\node[ecls, right=of source] (transcript)
  {Transcript};
\node[ecls, right=of transcript] (prompt)
  {Prompt\\[-2pt]\tiny promptHash};
\node[pkg, fit=(claim)(prompt)] (epkg) {};
\node[pkl, anchor=south west] at (epkg.north west)
  {Entities \(\sqsubseteq\) prov:Entity};
% Activities (middle)
\node[acls] (auth) at (0.6,-1.75) {AuthoringPass};
\node[acls, right=of auth] (audit) {AuditPass};
\node[acls, right=of audit] (build) {Build};
\node[acls, right=of build] (repair) {Repair};
\node[acls, right=of repair, draw=black!35, fill=black!5]
  (tele) {\tiny session, tokens,\\[-3pt]\tiny cost, energy, tools};
\node[pkg, fit=(auth)(tele)] (apkg) {};
\node[pkl, anchor=south west] at ([xshift=2mm]apkg.north west)
  {Activities \(\sqsubseteq\) prov:Activity};
% Agents (bottom)
\node[gcls] (human) at (1.4,-3.5) {HumanAgent\\[-2pt]\tiny orcid, affiliation};
\node[gcls, right=of human] (ai) {AIAgent\\[-2pt]\tiny model, provider};
\node[gcls, right=of ai] (toolag) {ToolAgent\\[-2pt]\tiny e.g.\ CI};
\node[pkg, fit=(human)(toolag)] (gpkg) {};
\node[pkl, anchor=south west] at ([xshift=2mm]gpkg.north west)
  {Agents \(\sqsubseteq\) prov:Agent};
% properties
\draw[e] ([xshift=-16mm]epkg.south) -- ([xshift=-19.5mm]apkg.north)
  node[lbl, midway] {wasGeneratedBy};
\draw[e] ([xshift=17mm]apkg.north) -- ([xshift=20mm]epkg.south)
  node[lbl, midway] {used, hadPlan, transcript};
\draw[e] (apkg.south) -- (gpkg.north)
  node[lbl, midway] {wasAssociatedWith};
\draw[e] (claim.west) to[bend right=40]
  node[lbl, pos=0.45, rotate=90] {wasAttributedTo} (gpkg.west);
\end{tikzpicture}%
}
\caption{Core aiprov vocabulary as an extension of PROV-O. Every class
subclasses a \texttt{prov:} term; tiny annotations name characteristic
attributes. Claims point to their generating activity, their
attributed agent, and a verification state; activities carry inference
telemetry and link prompts (as plans) and transcripts.}
\label{fig:vocab}
\end{figure}

\begin{table*}[t]
\caption{The verification ladder. Positions are machine-readable
(\texttt{aiprov:ladderPosition}); promotions must strictly increase
them and are themselves logged as audit activities, so every rung a
claim or source holds has a recorded grantor.}
\label{tab:ladder}
\centering
\small
\begin{tabular}{clp{8.2cm}l}
\toprule
Pos & Rung & What it asserts & Who may grant \\
\midrule
0 & \texttt{unverified} & Recorded; nothing checked yet by anyone.
    Default state at creation. & (default) \\
1 & \texttt{needs-research} & A check was attempted or demanded and did
    \emph{not} succeed (e.g.\ a DOI failed to resolve). Must not be
    cited or relied upon until re-verified. & any agent \\
2 & \texttt{reference-resolved} & The reference \emph{exists}: its
    DOI/URL resolved in a bibliographic registry and metadata was
    captured. Says nothing about content. & any agent (tool-checkable) \\
3 & \texttt{ai-confirmed} & An AI checked that the source
    \emph{content} supports the associated claim. A machine judgement,
    the highest rung an AI may grant. & AI or human \\
4 & \texttt{source-vendored} & A copy of the source is preserved inside
    the repository (content-hashed), immune to link rot and silent
    revision. Of no evidential value by itself: this rung is the
    \emph{access gate} to human verification: the agent hands the
    human the evidence, and tooling refuses rungs 5--6 unless the
    source is vendored or carries a clear access link (DOI/URL). &
    any agent (hash-checkable) \\
5 & \texttt{human-confirmed} & A \emph{human} spot-checked that the
    source supports the claim. & human only \\
6 & \texttt{human-read} & A \emph{human} has read the source in full
    \emph{and}, with that full context, confirms the claim is
    supported, subsuming rung 5's confirmation with deeper
    familiarity. Reading without confirming is not a rung. Top of the
    ladder. & human only \\
\bottomrule
\end{tabular}
\end{table*}

\begin{figure}[t]
\centering
\resizebox{0.80\columnwidth}{!}{%
\begin{tikzpicture}[
  rung/.style={draw, rounded corners=1pt, minimum width=5.6cm,
               minimum height=0.52cm, font=\footnotesize, align=center},
  human/.style={rung, fill=black!12},
  note/.style={font=\scriptsize\itshape, anchor=west}]
\node[human] (r6) {6\; human-read};
\node[human, below=0.16cm of r6] (r5) {5\; human-confirmed};
\node[rung, below=0.16cm of r5] (r4) {4\; source-vendored};
\node[rung, below=0.16cm of r4] (r3) {3\; ai-confirmed};
\node[rung, below=0.16cm of r3] (r2) {2\; reference-resolved};
\node[rung, below=0.16cm of r2] (r1) {1\; needs-research};
\node[rung, below=0.16cm of r1] (r0) {0\; unverified};
\node[note, right=0.15cm of r6] {human-only};
\node[note, right=0.15cm of r5] {human-only};
\node[note, right=0.15cm of r4] {access gate for 5--6};
\node[note, right=0.15cm of r3] {highest AI-grantable};
\draw[-{Stealth}, thick] ([xshift=-0.5cm]r0.west) -- ([xshift=-0.5cm]r6.west)
  node[midway, sloped, above, font=\scriptsize] {promotions strictly climb};
\end{tikzpicture}%
}
\caption{The verification ladder. Every promotion is logged as an audit
activity; rungs 5--6 may never be granted by an AI agent, and the
validator treats such grants as hard failures.}
\label{fig:ladder}
\end{figure}
